\documentclass{../../oss-apphys-exam}

\begin{document}
\genheader

\gentitle{C}{CIRCUIT ANALYSIS}

\genmultidirections

\gengravity

\raggedcolumns
\begin{multicols*}{2}

  \textbf{Questions \ref{circ1}--\ref{circ4}}  
  \begin{center}
    \begin{tikzpicture}[american voltages,scale=1.1]
      \draw(0,0)
      to[battery=\mbox{$\Delta V=\SI{12.}{\volt}$}] (0,3)
      to[R=\mbox{$R_1=\SI{3.}{\ohm}$}] (2,3)
      to[R=\mbox{$R_2=\SI{1.}{\ohm}$}] (2,0)
      to[R=\mbox{$R_3=\SI{2.}{\ohm}$}] (0,0);
    \end{tikzpicture}
  \end{center}
  
  \begin{questions}

    \question What is the current flowing through the circuit shown in the
    diagram?
    \label{circ1}
    \begin{choices}
      \choice\SI{1}{\ampere}
      \choice\SI{2}{\ampere}
      \choice\SI{4}{\ampere}
      \choice\SI{6}{\ampere}
      \choice\SI{12}{\ampere}
    \end{choices}
    
    \question Which of the following statements is true about the circuit shown
    in the diagram?
    \begin{choices}
      \choice The voltage drop is greatest across $R_1$, but $R_1$ has the least
      amount of current flowing through it.
      \choice The voltage drop is greatest across $R_2$, but $R_2$ has the least
      amount of current flowing through it.
      \choice The voltage drop is greatest across $R_3$, but $R_3$ has the least
      amount of current flowing through it.
      \choice The voltage drops and current are equal across all resistors.
      \choice The voltage drop is greatest across $R_1$, but the current is
      equal at all points in the circuit.
    \end{choices}
    
    \question In this diagram, what is the power dissipated by all of the
    resistors in the circuit?
    \begin{choices}
      \choice\SI{2}{\watt}
      \choice\SI{6}{\watt}
      \choice\SI{12}{\watt}
      \choice\SI{24}{\watt}
      \choice\SI{48}{\watt}
    \end{choices}
    \columnbreak
    
    \question In the diagram, what is the voltage drop across the third resistor
    ($R_3$)?
    \label{circ4}
    \begin{choices}
      \choice\SI{2}{\volt}
      \choice\SI{3}{\volt}
      \choice\SI{4}{\volt}
      \choice\SI{6}{\volt}
      \choice\SI{12}{\volt}
    \end{choices}

    \question Which of the following statements best summarizes a series circuit
    with three different resistances?
    \begin{choices}
      \choice In all parts of the circuit, the resistances are different, the
      voltage drops are the same, and the current is different.
      \choice In all parts of the circuit, the resistances are the same, the
      voltage drops are the same, and the current is different.
      \choice In all parts of the circuit, the resistances are different, the
      voltage drops are different, and the current is the same.
      \choice In all parts of the circuit, the resistances are different, the
      voltage drops are the same, and the current is the same.
      \choice In all parts of the circuit, the resistances are the same, the
      voltage drops are the same, and the current is the same.
    \end{choices}
    \columnbreak

    \uplevel{
      \textbf{Questions \ref{circ2-1}--\ref{circ2-2}}
      \begin{center}
        \begin{tikzpicture}[american voltages,scale=1.3]
          \draw(0,0) to[battery=$\Delta V$] (0,2)--(3,2)--(3,1.8);
          \draw(2,1.8)--(4.4,1.8);
          \draw(2,0.2)--(4.4,0.2);
          \draw(2,1.8) to[R=$R_1$] (2,0.2);
          \draw(3.2,1.8) to[R=$R_2$] (3.2,0.2);
          \draw(4.4,1.8) to[R=$R_3$] (4.4,0.2);
          \draw(3,0.2)--(3,0)--(0,0);
        \end{tikzpicture}
    
        \vspace{.1in}$\Delta V=\SI{12.0}{V}$, $R_1=\SI{10.0}{\ohm}$,\\
        $R_2=\SI{6.0}{\ohm}$, $R_3=\SI{8.0}{\ohm}$
      \end{center}
    }

    \question For the circuit in the diagram, which of the following
    expressions will describe the amount of current flowing through the
    resistors?
    \label{circ2-1}
    \begin{choices}
      \choice $I_1=I_2=I_3$
      \choice $I_3>I_2>I_1$
      \choice $I_1>I_2<I_3$
      \choice $I_2>I_1>I_3$
      \choice $I_1<I_2<I_3$
    \end{choices}
    
    \question For the circuit in the diagram, what is the equivalent resistance?
    \begin{choices}
      \choice\SI{0.040}{\ohm}
      \choice\SI{0.40}{\ohm}
      \choice\SI{1.0}{\ohm}
      \choice\SI{2.6}{\ohm}
      \choice\SI{24}{\ohm}
    \end{choices}
    
    \question For the circuit in the diagram, what is the total current?
    \begin{choices}
      \choice\SI{0.5}{\ampere}
      \choice\SI{4.6}{\ampere}
      \choice\SI{12}{\ampere}
      \choice\SI{30}{\ampere}
      \choice\SI{300}{\ampere}
    \end{choices}
    
    \question For the circuit in the diagram, the third resistor ($R_3$)
    dissipates how much energy each second?
    \label{circ2-2}
    \begin{choices}
      \choice\SI{12}{\watt}
      \choice\SI{14}{\watt}
      \choice\SI{46}{\watt}
      \choice\SI{212}{\watt}
      \choice\SI{300}{\watt}
    \end{choices}
    
%  \item Two resistors made of the same material are shown in the figure. A
%    current of $I$ flows through the left resistor when connected to a
%    potential difference of $V$. What current will flow through the right
%    resistor when connected to the same potential?
%
%    \vspace{-.2in}
%    \begin{center}
%      \pic{.45}{resistors.png}
%    \end{center}
%    \begin{enumerate}[noitemsep,topsep=0pt,leftmargin=18pt,label=(\Alph*)]
%    \item $I/2$
%    \item $I/4$
%    \item $I$
%    \item $2I$
%    \item $4I$
    %    \end{enumerate}
    \columnbreak
    
    \uplevel{
      \textbf{Questions \ref{series1}--\ref{series2}}
      \cpic{.45}{r-in-series}
    }

    \question Which is the correct ranking of the currents for the resistors?
    \label{series1}
    \begin{choices}
      \choice$I_A = I_B = I_C$
      \choice$I_A > I_B > I_C$
      \choice$I_C > I_A = I_B$
      \choice$I_C > I_B > I_A$
      \choice$I_C < I_B < I_A$
    \end{choices}

    \question Which is the correct ranking of the potential differences of the
    resistors?
    \label{series2}
    \begin{choices}
      \choice $V_A = V_B = V_C$
      \choice $V_A > V_B > V_C$
      \choice $V_A = V_B > V_C$
      \choice $V_C > V_B > V_A$
      \choice $V_C < V_B < V_A$
    \end{choices}

    \uplevel{
      \textbf{Questions \ref{combo1}--\ref{combo2}}
      \begin{center}
        \begin{tikzpicture}[american voltages,scale=1.2]
          \draw(0,0) to[battery,l=\mbox{\SI{120}{\volt}}] (0,2)
          to[C,l=\mbox{\SI{2}{\micro\farad}}] (2,2)
          to[C,l_=\mbox{\SI8{\micro\farad}}] (2,0)
          to[C,l=\mbox{\SI4{\micro\farad}}](0,0);
          \draw(2,2)--(3,2) to[C=\mbox{\SI8{\micro\farad}}] (3,0)--(2,0);
        \end{tikzpicture}
      \end{center}
    }

    \question The equivalent capacitance of this circuit is
    \label{combo1}
    \begin{choices}
      \choice\SI{7/4}{\micro\farad}
      \choice\SI{4/7}{\micro\farad}
      \choice\SI{21/16}{\micro\farad}
      \choice\SI{10}{\micro\farad}
      \choice\SI{22}{\micro\farad}
    \end{choices}
    
    \question The charge stored on the \SI{2}{\micro\farad} capacitor is most
    nearly
    \label{combo2}
    \begin{choices}
      \choice\SI{6}{\micro\coulomb}
      \choice\SI{12}{\micro\coulomb}
      \choice\SI{22}{\micro\coulomb}
      \choice\SI{36}{\micro\coulomb}
      \choice\SI{120}{\micro\coulomb}
    \end{choices}
    \columnbreak
    
    \question A capacitor $C_0$ is connected to a battery and stores charge. If
    the space between the capacitor plates is filled with oil, which of the
    following quantities increase?
    \begin{choices}
      \choice Capacitance and voltage across the plates
      \choice Charge and voltage across the plates
      \choice Capacitance and electric field between the plates
      \choice Capacitance and charge on the plates
      \choice Electric field between the plates and voltage across the plates
    \end{choices}

    \uplevel{
      \textbf{Question \ref{switch1}--\ref{switch2}}

      The circuit shows a capacitor, a battery, and a resistor. Switch $S$ is
      first connected to point $a$ to charge the capacitor, then a long time
      later switched to point $b$ to discharge the capacitor through the
      resistor.
      \begin{center}
        \begin{tikzpicture}[american voltages,scale=1.5]
          \draw(0,0) to[battery,l=$\mathcal{E}$](0,2) to[short,-*] (1.5,2);
          \draw(2.5,2) to[short,*-] (3.5,2) to[R=$R$,] (3.5,0)--(0,0);
          \draw(2,0) to[C=$C$,-*] (2,1.25);
          \draw(2,1.25)--(2.3,2.15) node[above]{s};
          \node at (1.5,1.8) (a){a};
          \node at (2.5,1.8) (b){b};
        \end{tikzpicture}
      \end{center}
    }

    \question The time constant $\tau$ for discharging the capacitor through the
    resistor could be decreased (faster discharge) by
    \label{switch1}
    \begin{choices}
      \choice placing another resistor in series with the first resistor
      \choice placing another resistor in parallel with the first resistor
      \choice placing another capacitor in parallel with the first capacitor
      \choice placing another battery in series in the same direction with the
      first battery
      \choice increasing both R and C
    \end{choices}
    
    \question The maximum current through the resistor is
    \label{switch2}
    \begin{choices}
      \choice $\mathcal{E}/2R$
      \choice $\mathcal{E}/R$
      \choice $\mathcal{E}/RC$
      \choice $\mathcal{E}/2RC$
      \choice $C\mathcal{E}/R$
    \end{choices}
    \columnbreak

    \question A battery of voltage $V_0$ is attached to two parallel conducting
    plates. Charge is distributed on the plates, and then the battery is
    removed. A dielectric is then inserted between the plates, filling the
    space. Which of the following decreases after the battery is removed and the
    dielectric is inserted to fill the space between the plates?
    \begin{choices}
    \item Capacitance
    \item Charge on the plates
    \item Net electric field between the plates
    \item Area of the plates
    \item Separation distance between the plates
    \end{choices}
    
    \question Circuit I and Circuit II shown each consist of a capacitor and a
    resistor. A battery is connected across a and b, and then removed. Which of
    the following statements is true of the circuits?
    \begin{center}
      Circuit I\\
      \begin{tikzpicture}
        \draw(0,1.25) to[short,*-](0,2.5)--(2,2.5) to[C=$C$](2,1.25)
        to[R=$R$](2,0)--(0,0) to[short,-*](0,0.75);
        \node at (-0.3,0.75) {a};
        \node at (-0.3,1.25) {b};
      \end{tikzpicture}
      
      Circuit II\\
      \begin{tikzpicture}
        \draw(0,1.25) to[short,*-](0,2)--(1.5,2) to[R=$R$](1.5,0)--(0,0)
        to[short,-*](0,0.75);
        \draw(1.5,2)--(3,2) to[C=$C$](3,0)--(1.5,0);
        \node at (-0.3,0.75) {a};
        \node at (-0.3,1.25) {b};
      \end{tikzpicture}
    \end{center}
    \begin{choices}
      \choice Circuit I and Circuit II will both retain stored energy when the
      battery is removed.
      \choice Neither Circuit I nor Circuit II will retain stored energy when
      the battery is removed.
      \choice Only Circuit I will retain stored energy when the battery is
      removed.
      \choice Only Circuit II will retain stored energy when the battery is
      removed.
      \choice Current will continue to flow in both circuits after the battery
      is removed.
    \end{choices}
    \columnbreak
    
    \question A thin sheet of copper is placed in a uniform magnetic field. A
    battery is connected to the top and bottom ends of the copper sheet, so that
    conventional current flows from the top to the bottom of the sheet.
    Points X and Y are on the left and right sides of the sheet, respectively.
    Which of the following statements is true?
    \begin{center}
      \begin{tikzpicture}[scale=.7]
        \foreach\x in {0,...,3}{
          \foreach \y in {0,...,3} \fill[gray](\x,\y) circle(.08); 
        }
        \draw[thick](.3,-.2) rectangle(2.7,3.2);
        \draw[thick](1.5,-.2)--(1.5,-.7)--(4.5,-.7)
        to[battery1,l_=$\mathcal{E}$](4.5,3.7)--(1.5,3.7)--(1.5,3.2);
        \fill(.3,1.5) circle(.1) node[left]{$X$};
        \fill(2.7,1.5) circle(.1) node[right]{$Y$};
        \node at (1,2.2) {$\mb{B}$};
      \end{tikzpicture}
    \end{center}
    \begin{choices}
      \choice Point $X$ is at a higher potential than point $Y$.
      \choice Point $Y$ is at a higher potential than point $X$.
      \choice Point $X$ and point $Y$ are at equal potential.
      \choice Point $X$ is at zero potential, and point $Y$ has a positive
      potential.
      \choice Point $Y$ is at zero potential, and point $X$ has a negative
      potential.
    \end{choices}

%  \item A very long conducting slab of copper of height $a$ and width $b$
%    carries a current $I$ throughout its cross-sectional area. The current
%    density $j$ is constant throughout the slab, and is directed out of the
%    page through the facing area of the slab. Points $X$ and $Y$ are marked on
%    the facing area of the slab. The current density $j$ can be expressed by
%    the expression
%    \begin{center}
%      \pic{.25}{square-conductor.png}
%    \end{center}
%    \begin{enumerate}[noitemsep,topsep=0pt,leftmargin=18pt,label=(\Alph*)]
%    \item $\displaystyle\frac{I}{a^2}$
%    \item $\displaystyle\frac{I}{b^2}$
%    \item $\displaystyle\frac{I}{(ab)^2}$
%    \item $\displaystyle\frac{I}{ab}$
%    \item $\displaystyle\frac{I}{a+b}$
%    \end{enumerate}
  \end{questions}
\end{multicols*}
\newpage

\genfreetitle{C}{ELECTRIC CIRCUITS}{5}

\genfreedirections

%TAKEN FROM 2010 AP PHYSICS EXAM FREE-RESPONSE QUESTION E&M 2
\cpic{.25}{r-c-1}

\begin{questions}
  \question In the circuit illustrated above, switch $S$ is initially open and
  the battery has been connected for a long time.
  \begin{parts}
    \part What is the steady-state current through the ammeter?
    
    \part Calculate the charge on the \SI{10}{\micro\farad} capacitor.
    
    \part Calculate the energy stored in the \SI{5.}{\micro\farad} capacitor.

    \uplevel{
      The switch is now closed, and the circuit comes to a new steady state.
    }

    \part Calculate the steady-state current through the battery.
    
    \part Calculate the final charge on the \SI{5.}{\micro\farad} capacitor.
    
    \part Calculate the energy dissipated as heat in the \SI{40}{\ohm} resistor
    in one minute once the circuit has reached steady state.
  \end{parts}
  \newpage

  
  % TAKEN FROM THE 2012 AP PHYSICS C EXAM FREE-RESPONSE QUESTION E&M 2
  \uplevel{
    \cpic{.5}{resistance}
  }
  \question A physics student wishes to measure the resistivity of slightly
  conductive paper that has a thickness of \SI{1.0e4}{\metre}. The student cuts
  a sheet of the conductive paper into strips of width \SI{.02}{\metre} and
  varying lengths, making five resistors labeled R1 to R5. Using an ohmmeter,
  the student measures the resistance of each strip, as shown above. The data
  are recorded below.
  \begin{center}
    \begin{tabular}{|l|c|c|c|c|c|}
      \hline
      Resistor & R1 & R2 & R3 & R4 & R5 \\ \hline
      Length (m) & 0.020 & 0.040 & 0.060 & 0.080 & 0.100\\ \hline
      Resistance (\si{\ohm}) & 80,000 & 180,000 & 260,000 & 370,000 & 440,000\\
      \hline
    \end{tabular}
  \end{center}
  \begin{parts}
    \part Use the grid below to plot a linear graph of the data points from
    which the resistivity of the paper can be determined. Include labels and
    scales for both axes. Draw the straight line that best represents the data.
    \begin{center}
      \begin{tikzpicture}[yscale=.3,xscale=.45]
        \draw (0,0) grid(30,30);
        \draw[step=5,very thick](0,0) grid(30,30);
      \end{tikzpicture}
    \end{center}

    \part Using the graph, calculate the resistivity of the paper.

    \uplevel{
      \cpic{.45}{resistors-and-capacitor}
      The student uses resistors R4 and R5 to build a circuit using wire, a
      \SI{1.5}{\volt} battery, an uncharged \SI{10}{\micro\farad} capacitor,
      and an open switch, as shown above.
    }
    
    \part Calculate the time constant of the circuit.
    \newpage
    
    \part At time $t=0$, the student closes the switch. On the axes below,
    sketch the magnitude of the voltage $V_c$ across the capacitor and the
    magnitudes of the voltages $V_\text{R4}$ and $V_\text{R5}$ across each
    resistor as functions of time $t$. Clearly label each curve according to
    the circuit element it represents. On the axes, explicitly label any
    intercepts, asymptotes, maxima, or minima with values or expressions, as
    appropriate.
    \begin{center}
      \begin{tikzpicture}[scale=.9]
        \draw[very thick,->](0,0)--(10,0)
        node[midway,below]{Time $t$} node[pos=0,below left]{$O$};
        \draw[very thick,->](0,0)--(0,6)node[midway,above,sloped]{Voltage $V$};
      \end{tikzpicture}
    \end{center}
  \end{parts}
  \newpage
  
%\item The circuit in the figure consists of three capacitors
%  (\SI{3}{\micro\farad}, \SI{4}{\micro\farad}, and \SI{6}{\micro\farad})
%  connected to a \SI{200}{\volt} battery.
%  
%  \begin{minipage}{0.3\textwidth}
%    \begin{tikzpicture}[american voltages,scale=1.3]
%      \draw[thick](4,0) to[battery=200<\volt>](0,0)--(0,1)
%      to[C=4<\micro\farad>](4,1)--(4,0);
%      \draw[thick](0,1)--(0,2) to[C=6<\micro\farad>] (2,2)
%      to[C=3<\micro\farad>] (4,2)--(4,1);
%    \end{tikzpicture}
%  \end{minipage}
%  \begin{minipage}{.68\textwidth}
%    \begin{enumerate}[noitemsep]
%    \item Calculate the equivalent capacitance of the combined three
%      capacitors.
%    \item Calculate the total energy stored in the \SI{6}{\micro\farad} and
%      \SI{3}{\micro\farad} capacitor combination.
%    \end{enumerate}
%  \end{minipage}


%TAKEN FROM 2007 AP PHYSICS C EXAM FREE-RESPONSE QUESTION E&M 1.
  \uplevel{
    \cpic{.4}{r-c-2}
  }
  \question A student sets up the circuit above in the lab. The values of the
  resistance and capacitance are as shown, but the constant voltage
  $\mathcal{E}$ delivered by the ideal battery is unknown. At time $t=0$, the
  capacitor is uncharged and the student closes the switch. The current as a
  function of time is measured using a computer system, and the following graph
  is obtained.
  \cpic{.4}{current-graph}
  \begin{parts}
    \part Using the data above, calculate the battery voltage $\mathcal{E}$.

    \part Calculate the voltage across the capacitor at time
    $t =\SI{4.}{\second}$.
    
    \part Calculate the charge on the capacitor at $t =\SI{4.}{\second}$.

    \part On the axes below, sketch a graph of the charge on the capacitor as a
    function of time.
    \begin{center}
      \begin{tikzpicture}[scale=.65]
        \draw[very thick,->](0,0)--(10,0)
        node[midway,below]{Time ($s$)} node[pos=0,below left]{$O$};
        \draw[very thick,->](0,0)--(0,6)node[midway,above,sloped]{Charge};
        \draw(9,-.1)--(9,.1) node[pos=0,below]{10};
      \end{tikzpicture}
    \end{center}
    \part Calculate the power being dissipated as heat in the resistor at
    $t=\SI{4.}{\second}$.

    \part  The capacitor is now discharged, its dielectric of constant
    $\kappa=1$ is replaced by a dielectric of constant $\kappa=3$, and the
    procedure is repeated. Is the amount of charge on one plate of the
    capacitor at $t=\SI{4.}{\second}$ now greater than, less than, or the same
    as before? Justify your answer.

    \vspace{.1in}
    \underline{\hspace{.3in}} Greater than\hspace{.3in}
    \underline{\hspace{.3in}} Less than\hspace{.3in}
    \underline{\hspace{.3in}} The same
  \end{parts}
  \newpage
  
  %TAKEN FROM 2017 AP PHYSICS C EXAM FREE-RESPONSE QUESTION E&M 2
  \uplevel{
    \cpic{.4}{r-c-3}
  }
  \question In the circuit above, an ideal battery of voltage $V_0$ is
  connected to a capacitor with capacitance $C_0$ and resistors with
  resistances $R_1$ and $R_2$ , with $R_1>R_2$. The switch $S$ is open, and the
  capacitor is initially uncharged.
  \begin{parts}
    \part The switch is closed at time $t=0$. On the axes below, sketch the
    charge $q$ on the capacitor as a function of time $t$. Explicitly label any
    intercepts, asymptotes, maxima, or minima with numerical values or algebraic
    expressions, as appropriate.
    \begin{center}
      \begin{tikzpicture}[scale=.65]
        \draw[very thick,->](0,0)--(10,0) node[right]{$t$};
        \draw[very thick,->](0,0)--(0,5)  node[above]{$q$};
      \end{tikzpicture}
    \end{center}

    \part On the axes below, sketch the current $I$ through each resistor as a
    function of time $t$. Clearly label the two curves as $I_1$ and $I_2$ , the
    currents through resistors $R_1$ and $R_2$, respectively. Explicitly label
    any intercepts, asymptotes, maxima, or minima with numerical values or
    algebraic expressions, as appropriate.
    \begin{center}
      \begin{tikzpicture}[scale=.65]
        \draw[very thick,->](0,0)--(10,0) node[right]{$t$};
        \draw[very thick,->](0,0)--(0,5)  node[above]{$I$};
      \end{tikzpicture}
    \end{center}

    \uplevel{
      The circuit is constructed using an ideal \SI{1.5}{\volt} battery, an
      \SI{80}{\micro\farad} capacitor, and resistors $R_1=\SI{150}{\ohm}$ and
      $R_2=\SI{100}{\ohm}$. The switch is closed, allowing the capacitor to
      fully charge. The switch is then opened, allowing the capacitor to
      discharge.
    }

    \part The time it takes to charge the capacitor to \SI{50}{\percent} of its
    maximum charge is $\Delta t_C$. The time it takes for the capacitor to
    discharge to \SI{50}{\percent} of its maximum charge is $\Delta t_D$.
    Which of the following correctly relates the two time intervals?
    Justify your answer.
    
    \vspace{.1in}
    \underline{\hspace{.3in}} $\Delta t_C > \Delta t_D$\hspace{.3in}
    \underline{\hspace{.3in}} $\Delta t_C = \Delta t_D$\hspace{.3in}
    \underline{\hspace{.3in}} $\Delta t_C < \Delta t_D$
    
    \part
    \begin{subparts}
      \subpart Calculate the current through resistor $R_2$ immediately after
      the switch is opened.
      \subpart Is the current through resistor $R_2$ increasing, decreasing, or
      constant immediately after the switch is opened?
      Justify your answer.
      
      \vspace{.1in}
      \underline{\hspace{.3in}} Increasing\hspace{.3in}
      \underline{\hspace{.3in}} Decreasing\hspace{.3in}
      \underline{\hspace{.3in}} Constant
    \end{subparts}
    \part
    \begin{subparts}
      \subpart Calculate the energy stored in the capacitor immediately after
      the switch is opened.
      \subpart Calculate the energy dissipated by resistor $R_1$ as the
      capacitor completely discharges.
    \end{subparts}
  \end{parts}
  \newpage
  
  %TAKEN FROM THE 2001 AP PHYSICS C FREE-RESPONSE QUESTION E&M 2

  \question You have been hired to determine the internal resistance of
  \SI{8.}{\micro\farad} capacitors for an electronic component manufacturer.
  (Ideal capacitors have an infinite internal resistance—that is, the material
  between their plates is a perfect insulator. In practice, however, the
  material has a very small, but nonzero, conductivity.) You cannot simply
  connect the capacitors to an ohmmeter, because their resistance is too large
  for an ohmmeter to measure. Therefore you charge the capacitor to a potential
  difference of \SI{10}{\volt} with a battery, disconnect it from the battery
  and measure the potential difference across the capacitor every 20 minutes
  with an ideal voltmeter, obtaining the graph shown above.
  \begin{parts}
    \part Determine the internal resistance of the capacitor.
    
    \uplevel{
      The capacitor can be approximated as a parallel-plate capacitor separated
      by a \SI{.10}{\milli\metre} thick dielectric with $\kappa=5.6$.
    }

    \part Determine the approximate surface area of one of the capacitor
    ``plates.''
    
    \part Determine the resistivity of the dielectric.
    
    \part Determine the magnitude of the charge leaving the positive plate of
    the capacitor in the first \SI{100}{min}.
  \end{parts}
  
%\item A \SI{60}{\micro\farad} capacitor is charged to \SI{12}{\volt}. The
%  capacitor is then removed from the battery and the plate separation is
%  increased from \SI{2.0}{mm} to \SI{3.5}{mm}.
%  \begin{enumerate}[noitemsep]
%  \item What is the charge on the capacitor?
%  \item How much energy was originally stored in the capacitor?
%  \item By how much is the energy increased when the plate separation is
%    changed?
%  \end{enumerate}
%  \newpage

%\item Find the current in each part of the circuit shown in the figure below.
%  \begin{center}
%    \begin{tikzpicture}[american voltages,scale=1.25]
%      \draw[thick](0,0) to[short,*-](0,2) to[R=3<\ohm>,*-] (0,4)
%      to[battery1=42<\volt>,*-](2,4) to[R=3<\ohm>] (4,4) to[short,*-](4,2)
%      to[R=4<\ohm>,*-](4,0) to[short,*-](2,0) to[battery1=6<\volt>,*-](2,2)
%      to[R,l_=4<\ohm>,*-](0,2);
%      \draw[thick](0,0) to[R,l_=6<\ohm>](2,0);
%      \draw[thick](2,2) to[R,l_=6<\ohm>](4,2);
%    \end{tikzpicture}
%  \end{center}
%  \newpage
%  
%\item The capacitor in the circuit is initially uncharged. Find the current
%  through the battery
%  \begin{enumerate}[noitemsep]
%  \item Immediately after the switch is closed
%  \item A long time after the switch is closed 
%  \end{enumerate}
%
%  \begin{center}
%    \begin{tikzpicture}[scale=1.25,american voltages]
%      \draw[thick] (0,0) to[battery1=12<\volt>,*-*] (0,2) to[R=4<\ohm>] (2,2)
%      to[short,*-](2.46,2.3);
%      \draw[thick](2.5,2) to[short,*-*] (3,2) to[short,-*] (4,2)
%      to[C=6<\micro\farad>,-*] (4,0)--(0,0);
%      \draw[thick] (3,0) to[R=8<\ohm>,*-] (3,2);
%    \end{tikzpicture}
%  \end{center}
%
%
%\item The figure shows a circuit with two resistors, a battery, a capacitor,
%  and a switch. Originally, the switch is open, and the capacitor is
%  uncharged.
%
%  \vspace{-.5cm}
%  \begin{center}
%    \begin{tikzpicture}[scale=1.4,american voltages]
%      \draw[thick] (2,0) to[battery1=12<\volt>,*-*] (0,0)
%      to[short,-*] (0,1.5) to[R=\mbox{$R_1=\SI{15}{\ohm}$},-*] (2,1.5)--(2,0);
%      \draw[thick](0,1.5) to[short,-*] (0,3) to[C=240<\micro\farad>] (1,3)
%      to[R=\mbox{$R_2=\SI{10}{\ohm}$},-*] (2,3)--(2,1.5);
%    \end{tikzpicture}
%  \end{center}
%  
%  \begin{enumerate}[noitemsep]
%  \item Complete the voltage-current-resistance-power (VIRP) chart for
%    the circuit immediately after the switch is closed.
%    \bgroup
%    \def\arraystretch{2}
%    
%    {\large
%      \begin{center}
%        \begin{tabular}{l|m{1.2cm}|m{1.2cm}|m{1.2cm}|m{1.2cm}}
%          \hline
%          \textbf{Location} & \textbf{\emph{V}} & \textbf{\emph{I}} &
%          \textbf{\emph{R}} & \textbf{\emph{P}}\\ \hline
%          1 & & & \SI{15}{\ohm} & \\ \hline
%          2 & & & \SI{10}{\ohm} & \\ \hline
%          Total for Circuit & \SI{12}{\volt}& & & \\\hline
%        \end{tabular}
%      \end{center}
%    }
%    \egroup
%  \item Complete the voltage-current-resistance-power (VIRP) chart for
%    the circuit after the switch is closed for a long time.
%        \bgroup
%    \def\arraystretch{2}
%    
%    {\large
%      \begin{center}
%        \begin{tabular}{l|m{1.2cm}|m{1.2cm}|m{1.2cm}|m{1.2cm}}
%          \hline
%          \textbf{Location} & \textbf{\emph{V}} & \textbf{\emph{I}} &
%          \textbf{\emph{R}} & \textbf{\emph{P}}\\ \hline
%          1 & & & \SI{15}{\ohm} & \\ \hline
%          2 & & & \SI{10}{\ohm} & \\ \hline
%          Total for Circuit & \SI{12}{\volt}& & & \\\hline
%        \end{tabular}
%      \end{center}
%    }
%    \egroup
%  \item What is the energy stored in the capacitor after the switch has
%    been closed a long time?
%  \end{enumerate}
\end{questions}
\end{document}
